\section{STAKEHOLDER REQUIREMENTS}

The stakeholders for the Tri-Mic are:

\begin{enumerate}
	\item End users (field technicians and hardware enthusiasts)
	\item Environmental and durability concerns
	\item Data and export needs
	\item Manufacturing and engineering team
\end{enumerate}

\subsection{Stakeholder User Stories}

\begin{itemize}
	\item \textbf{The end user operates the Tri-Mic in the field.  The device:}
	\begin{enumerate}
		\item Must pinpoint specific sounds in 3D and isolate them from background noise.\\ [.5cm] \underline{User Story}\\
		I work around loud equipment and need to figure out which direction a problem sound is coming from.  Right now I just use my ears and walk around until I think I found it.  The mic should tell me the direction so I'm not guessing.
		\bigskip

		\item Must show a spectrogram on the device without needing a laptop.\\ [.5cm] \underline{User Story}\\
		I don't want to carry a laptop out to every job.  If I can see the frequency data right on the device I can make decisions on site and keep working.
		\bigskip

		\item Must let the user steer the directional focus and see the spectrogram change in real time.\\ [.5cm] \underline{User Story}\\
		I want a knob or joystick that I can move around and watch the display update as I sweep different directions.  Something physical makes the most sense---I might have gloves on or be in a tight spot where a touchscreen isn't going to work.
		\bigskip

		\item Must be small enough to fit into tight spaces.\\ [.5cm] \underline{User Story}\\
		Half the places I need to check are behind panels or inside enclosures.  The mic needs to get in there without me having to shove my head in.
		\bigskip
	\end{enumerate}

	\item \textbf{The device must handle real-world field conditions:}
	\begin{enumerate}
		\item Must be tough enough to handle drops and field use.\\ [.5cm] \underline{User Story}\\
		Stuff gets dropped.  If this breaks from falling off a ladder it's not going to work for us.
		\bigskip

		\item Must still work in damp conditions.\\ [.5cm] \underline{User Story}\\
		I work in crawl spaces, mechanical rooms, and sometimes outside.  A little moisture or humidity shouldn't kill it.
		\bigskip
	\end{enumerate}

	\item \textbf{Captured data needs to be useful after the fact:}
	\begin{enumerate}
		\item Must be able to export audio and directional data to a PC.\\ [.5cm] \underline{User Story}\\
		The on-device display is fine for field work but for real analysis we need to pull the data off and run it through our own tools on a computer.  We want WAV files and TDOA data we can actually work with.
		\bigskip
	\end{enumerate}

	\item \textbf{Manufacturing and engineering cares about cost and flexibility:}
	\begin{enumerate}
		\item Must use standard off-the-shelf components.\\ [.5cm] \underline{User Story}\\
		We don't want exotic parts with long lead times.  Standard mics, standard PCB components, standard connectors.  The MUSIC algorithm is well-established so we're not reinventing the wheel on the computation side either.
		\bigskip

		\item Must support adjustable mic spacing and swappable DSP pipelines during prototyping.\\ [.5cm] \underline{User Story}\\
		We need to try different array configurations to see what gives the best direction-finding accuracy.  MUSIC is our starting algorithm but we might try other approaches.  The architecture should let us change both the mic geometry and the processing without redesigning the hardware.
		\bigskip
	\end{enumerate}
\end{itemize}
